\section{评测污染与不可比设定审计}\label{sec:audit}

本文前半部分引用了大量文献结论。若不加审计地引用，就会把一个已知的问题带入结论：不同论文的绩效数字建立在不可比的设定上，且普遍存在数据窥探。本章按八个核验维度做审计，并\emph{同时审计本文自建回测的局限}——后者与前者同样重要，因为自建实证最容易被误当作无偏的裁判。

\subsection{审计框架}

\begin{table}[htbp]
\centering\footnotesize
\caption{八维核验口径在策略比较语境下的映射}\label{tab:audit}
\begin{tabular}{@{}>{\raggedright\arraybackslash}p{2.6cm}>{\raggedright\arraybackslash}p{11.2cm}@{}}
\toprule
\textbf{核验维度} & \textbf{本语境下的核验问题} \\
\midrule
样本数据 & 样本期、资产池、频率是否可比？是否含幸存者偏差与上市偏差？ \\
信号设定 & 形成期/持有期、阈值、参数是在样本内选的还是外生的？ \\
样本外稳健性 & 是否做 walk-forward 与跨市场泛化检验？ \\
概率/状态预测 & 是否给出参数估计与状态判断（OU 参数、半衰期、regime）？ \\
计算成本 & 数据与算力需求；重估频率 \\
许可 & 数据（CRSP/TAQ/交易所）与代码许可是否可得？ \\
公开基准 & 是否提供可复现基准（代码、复现报告）？ \\
数据污染与不可比 & look-ahead、数据窥探、多重检验、跨论文设定差异 \\
\bottomrule
\end{tabular}
\end{table}

\subsection{幸存者偏差与上市偏差}

两条脉络在这一项上的暴露程度不同。

趋势研究的资产池通常是\emph{存续至今}的期货合约或指数成分。Hurst 等(2017) 手工收集 1880 年以来的 CBOT 数据，所用 67 个市场是“可回溯到”的市场\cite{hop2017}；那些中途退市、停止交易的市场未被纳入。这构成典型的幸存者偏差：策略在“活下来的市场”上表现好，可能是选择的结果。

反转研究的偏差更隐蔽。配对交易需要在形成期就能配对，而“能配对”本身隐含了样本的存续要求。Gatev 等(2006) 用 1962--2002 年 CRSP 数据，其配对的股票池是存续股票\cite{gatev2006}；Do \& Faff(2012) 延伸样本至 2009 年后利润衰减，部分原因可能就是早期样本的存续偏差在后期被稀释\cite{dofaff2012}。

此外，反转策略还多一层偏差：\emph{它需要做空}。在多数市场中，可做空的股票集合远小于全部股票，且做空成本与可借量随时间变化。文献中的反转收益大多在“可做空”的假设下计算，实际可实现的收益更低。这一项在第~\ref{sec:empirical}~节的 A 股部分尤其致命，见后文审计。

\subsection{Look-ahead 偏差}

趋势研究的典型 look-ahead 来源是\emph{信号与执行的时间错配}：若在信号形成期的最后一个交易日收盘价上建仓，而实际执行发生在次日开盘，则回测高估收益。MOP12 的构造以月末信号、次月持有处理了这一项\cite{mop2012}。

反转研究的典型来源是\emph{配对选择与交易期重叠}。若配对是在包含交易期的数据上选出的，则价差平稳性被机械地保证。标准做法是形成期与交易期严格分离并滚动（Gatev 等的 6+6 个月滚动设计\cite{gatev2006}）。

本文自建回测的处理是：信号使用至 $t$ 日收盘的全部可得信息，仓位自 $t+1$ 日起生效，持有期内按收益漂移，调仓以目标权重与实际权重的偏离计换手。这一处理避免了最常见的 look-ahead，但仍有残余问题（见~\ref{sec:selfaudit}~小节）。

\subsection{数据窥探与多重检验}

这是两条脉络共同的、也是最严重的威胁。

\textbf{趋势一侧的历史更早被审查}。Sullivan 等(1999) 对 7846 条技术交易规则做 bootstrap 校正，发现样本外超额收益不再显著\cite{sullivan1999}；White(2000) 的 Reality Check 给出通用框架\cite{white2000}；Hansen(2005) 的 SPA 检验与 Romano--Wolf(2005) 的逐步多重检验进一步改进\cite{hansen2005,romanowolf2005}。Harvey 等(2016) 提出的更高门槛（新因子应达到 $t>3$ 左右）\cite{hlz2016}、以及 Bailey 与 López de Prado 的紧缩夏普比率与回测过拟合概率\cite{bailey2014,bailey2016}，构成当前的审计标准。

\textbf{反转一侧的暴露更隐蔽}。配对交易的超参数包括选对规则、协整显著性、半衰期上下限、阈值带宽度、形成期长度；每一个都是自由度。Krauss(2017) 的综述在梳理该领域时明确指出方法选择的空间巨大\cite{krauss2017}。这意味着反转文献中“发现显著收益”的比例可能被系统性高估，而 Do \& Faff(2012) 的利润衰减正是这一高估在样本外被修正的表征\cite{dofaff2012}。

\textbf{本文自建回测的暴露}。本文对 $L\in\{5,10,20,60,120,250\}$ 共六个窗口、两个方向、两种口径（时序/截面）、两个市场做了完整网格扫描，未做多重检验校正。因此本文报告的“某组合夏普为正”不应当被读作统计显著，只能读作\emph{在给定口径下的相对比较}。这是本文实证的定位：它是变量对照，不是显著性检验。

\subsection{样本外与发表后衰减}

McLean \& Pontiff(2016) 对 97 个已发表因子的实测是这一项最有力的证据：组合收益在样本外-发表前下降 26\%，发表后累计下降约 58\%\cite{mcleanpontiff2016}。这一结论对两条脉络同等适用——趋势与反转都是“被发表过的异象”。

对趋势策略的特别审查来自 Huang 等(2020)：逐资产回归显示时间序列动量的证据在样本内外均薄弱，池回归的大 $t$ 值不稳健，且策略表现与“按历史样本均值投资”的简单规则几乎相同\cite{huang2020}。Pitkäjärvi 等(2020) 从跨资产信号角度给出相关质疑\cite{pitkajarvi2020}。反方证据是 Hurst 等(2017) 的百年样本\cite{hop2017}——两方证据的张力在结论置信度章（第~\ref{sec:confidence}~节）中作为专门条目处理。

对反转策略的特别审查是利润的时间衰减：Gatev 等(2006) 报告的约 11\% 年化超额收益基于 1962--2002 样本\cite{gatev2006}，Do \& Faff(2012) 显示 1989 年之后明显下降\cite{dofaff2012}。这提示反转文献的高收益不能外推。

\subsection{跨论文不可比设定}

即便不考虑偏差，跨文献的绩效数字也常常不可比。差异来源至少四类：

\begin{enumerate}[leftmargin=2.2em]
  \item \textbf{样本期与资产池}：MOP12 用 1960s--2009 的 58 个合约\cite{mop2012}，HOP17 用 1880--2016 的 67 个市场\cite{hop2017}，Georgopoulou \& Wang(2017) 用 1969--2015 的股票与商品\cite{georgopoulou2017}；三者的“趋势跟踪年化”不是同一个量。
  \item \textbf{波动率缩放是否启用}：MOP12 的原始设定与 Barroso--Santa-Clara(2015)、Harvey 等(2018) 的缩放设定在风险暴露上不同\cite{mop2012,barroso2015,harvey2018}；Kim 等(2016) 指出收益的很大部分来自缩放本身\cite{kim2016}。
  \item \textbf{成本与费用假设}：HOP17 计入滚动成本与管理费\cite{hop2017}，部分复现研究忽略；反转一侧的成本假设差异更大（价差、融券费、冲击模型各不相同）\cite{dofaff2012,novymarxvelikov2016}。
  \item \textbf{频率}：月频、周频、日频策略的换手相差一个量级，成本敏感度随之变化\cite{baltaskosowski2011,korajczyksadka2004}。
\end{enumerate}

因此本文在引用任何绩效数字时均注明其设定条件，并在比较矩阵中拒绝把不同设定的数字直接排序。

\subsection{本文自建回测的局限}\label{sec:selfaudit}

自建实证同样需要审计，且这一小节是本文最应当被严格对待的部分。

\textbf{局限一：A 股个股采用不复权价的近似修复，存在残余污染}。通达信本地日线为不复权数据，本文的修复规则是——主板个股单日收益低于 $-10.5\%$、创业板/科创板低于 $-20.5\%$ 者判定为除权并置该日收益为零（等价于按当日比例对更早价格做前复权）。这一处理能捕获绝大多数送转除权（跳空幅度大于涨跌停），但\emph{漏掉跳空幅度小于涨跌停的除权}（如小额分红、10 送 1），后者构成残余污染。样本内共识别并修复 6121 处确定性除权；未识别的部分未做定量估计，属于本文已知的未修正偏差。

\textbf{局限二：A 股截面策略的做空假设不现实}。A 股个股的融券做空渠道有限、成本高、可借券池小。本文的 A 股截面多空组合因此在多数时段不可实现，其结果应被视为\emph{上限估计}而非可交易绩效。这一局限对反转策略的影响大于趋势策略，因为反转的多空结构更依赖空头端。

\textbf{局限三：未建模涨跌停与停牌的可交易性}。涨跌停板上无法成交，停牌期间无法调仓。本文的信号与收益计算未做这一过滤，因此在高波动时段（如 2015 年 A 股异常波动期间）会高估策略的可实现收益。期货侧的涨跌停同样未过滤。

\textbf{局限四：成本模型是简化的固定基点}。本文对期货采用单边 2 个基点、A 股采用单边 15 个基点（含佣金、印花税与滑点），未采用参与率约束的平方根冲击模型，也未考虑流动性时变。这一简化会低估反转策略（高换手）的成本，因为反转交易密集发生在流动性较差的时点。

\textbf{局限五：日线口径}。所有策略在日线频率上构建，无法刻画日内反转与执行细节；这与反转策略的实际高频属性存在口径落差。

\textbf{局限六：样本期较短且单一市场}。2014--2026 的十二年间，中国期货与 A 股的制度环境经历了显著变化（股指期货限制、融券政策调整、注册制改革）；结论不应外推到更早时期或其他市场。

这些局限不改变本文实证的定位——它是一次\emph{变量受控的相对比较}（同一资产池、同一成本、仅翻转符号），而不是对绝对收益水平的估计。任何把本文数字当作“策略年化收益”的读法都是误读。

\subsection{审计小结}

两条脉络的污染暴露在结构上同构：都有数据窥探与发表后衰减，都有不可比的样本设定，都有做空与流动性的实现障碍。差异在于侧重点——趋势研究更早、更彻底地接受了多重检验审查；反转研究的自由度（配对选择、参数筛选）更多，因此更依赖于“利润衰减”这类事后证据来暴露问题。

对本文自建实证的审计结论是：它可以支撑\emph{方向性比较}（两策略在同一口径下的相对表现），不能支撑\emph{水平性结论}（某策略能赚多少）。
